\chapter{Detailed Derivations}
\label{app:derivations}

\section{Taylor expansion of binary entropy}
Let $\sigma=1/2+\delta$. Then
\begin{align}
\Hbin\!\left(\frac12+\delta\right)
={}&-\left(\frac12+\delta\right)
\ln\left(\frac12+\delta\right)
-\left(\frac12-\delta\right)
\ln\left(\frac12-\delta\right).
\end{align}
Using
\begin{equation}
\ln\left(\frac12\pm\delta\right)
=-\ln2+\ln(1\pm2\delta)
\end{equation}
and the power series for $\ln(1+x)$ gives
\begin{equation}
\Hbin\!\left(\frac12+\delta\right)
=\ln2-2\delta^2-\frac43\delta^4-\frac{32}{15}\delta^6+O(\delta^8).
\end{equation}
All odd powers cancel by complement symmetry.

\section{Taylor expansion of the cancellation residual}
At phase $\pi$,
\begin{equation}
\mathcal I
=\left(\sqrt{\frac12+\delta}-\sqrt{\frac12-\delta}\right)^2.
\end{equation}
Since
\begin{equation}
\mathcal I=1-2\sqrt{\frac14-\delta^2}
=1-\sqrt{1-4\delta^2},
\end{equation}
the binomial series yields
\begin{equation}
\mathcal I
=2\delta^2+2\delta^4+4\delta^6+10\delta^8+O(\delta^{10}).
\end{equation}
Thus the information deficit and interference residual coincide only to leading nonzero order; their higher-order coefficients differ.

\section{Fisher coordinate and Bloch angle}
For $\sigma=\cos^2(\theta/2)$,
\begin{equation}
\frac{\dd\sigma^2}{\sigma(1-\sigma)}
=\frac{(\frac14\sin^2\theta\,\dd\theta^2)}{\frac14\sin^2\theta}
=\dd\theta^2.
\end{equation}
The Fisher distance from the balanced point is therefore
\begin{equation}
d_F(\sigma,1/2)
=\left|\theta-\frac\pi2\right|
=\left|2\arccos\sqrt\sigma-\frac\pi2\right|.
\end{equation}

\section{Berry connection in two gauges}
The north-patch gauge used in the main text gives
\begin{equation}
\ket{\psi_N}
=\cos(\theta/2)\ket0+e^{i\phi}\sin(\theta/2)\ket1,
\end{equation}
with
\begin{equation}
\mathcal A_N=-\sin^2(\theta/2)\,\dd\phi
=-\frac{1-\cos\theta}{2}\dd\phi.
\end{equation}
A south-patch gauge may be chosen as
\begin{equation}
\ket{\psi_S}=e^{-i\phi}\cos(\theta/2)\ket0+\sin(\theta/2)\ket1,
\end{equation}
with
\begin{equation}
\mathcal A_S=\cos^2(\theta/2)\,\dd\phi
=\frac{1+\cos\theta}{2}\dd\phi.
\end{equation}
Their difference is
\begin{equation}
\mathcal A_S-\mathcal A_N=\dd\phi,
\end{equation}
so the integrated phases differ by $2\pi$ and the holonomy is identical.

\section{Modulus of the functional-equation factor}
On the critical line, the functional equation and conjugation give
\begin{equation}
\zeta(s)=\chi(s)\overline{\zeta(s)}.
\end{equation}
At a point where $\zeta(s)\neq0$, taking absolute values yields
\begin{equation}
|\zeta(s)|=|\chi(s)|\,|\zeta(s)|,
\end{equation}
so $|\chi(s)|=1$. The identity extends continuously through zeros. This argument uses the fixed-line relation $1-s=\overline s$ and therefore does not hold off the line.

\section{Half-density term in the dilation generator}
Define the unitary dilation group on $L^2(\mathbb R_+,\dd x)$ by
\begin{equation}
(U_a f)(x)=e^{a/2}f(e^a x).
\end{equation}
Unitarity follows from
\begin{align}
\|U_a f\|^2
&=\int_0^\infty e^a|f(e^a x)|^2\,\dd x\\
&=\int_0^\infty |f(y)|^2\,\dd y.
\end{align}
Differentiating at $a=0$ gives
\begin{equation}
\left.\frac{\dd}{\dd a}U_a f\right|_{a=0}
=\frac12 f+x f'.
\end{equation}
Therefore the self-adjoint generator convention $U_a=e^{-ia\widehat D}$ gives
\begin{equation}
\widehat D=-i\left(x\frac{\dd}{\dd x}+\frac12\right).
\end{equation}
The half term is required by the Jacobian of the unitary action.
